\begin{exercisechunk}
\item \textbf{Initialization and the minimum-norm least-squares solution.}
\label[exercise]{ex:l10-least-squares-selection}
\exerciseprofile{2}{2}{2}{\exproof\quad\excalculation}{%
Identify the invariant component of least-squares gradient descent and derive
the solution selected by its initialization.}
Fix $x_1,\ldots,x_n\in\R^d$ and $y_1,\ldots,y_n\in\R$. Let
\[
 \cR=\operatorname{span}\{x_1,\ldots,x_n\},
 \qquad \dim(\cR)=n,
\]
let $\mathbf X\in\R^{n\times d}$ have row $i$ equal to $x_i^\top$, and set
$y=(y_1,\ldots,y_n)^\top$. Define
\[
 L_n(w)=\frac12\sum_{i=1}^n
 \bigl(\langle w,x_i\rangle-y_i\bigr)^2.
\]
Consider gradient descent
\[
 w_{t+1}=w_t-
 \eta\sum_{i=1}^n
 \bigl(\langle w_t,x_i\rangle-y_i\bigr)x_i,
 \qquad 0<\eta<2/\|\mathbf X\|_{\mathrm{op}}^2.
\]
\begin{enumerate}[(a),beginpenalty=10000]
\item Derive the normal equation in both forms
\[
 \sum_{i=1}^n
 \bigl(\langle w,x_i\rangle-y_i\bigr)x_i=0,
 \qquad
 \mathbf X^\top(\mathbf Xw-y)=0.
\]
Prove that $\dim(\cR)=n$ is equivalent to
$\mathbf X\mathbf X^\top$ being nonsingular, and deduce that the normal
equation is equivalent to $\mathbf Xw=y$.

\item Define
$w^\star=\mathbf X^\top(\mathbf X\mathbf X^\top)^{-1}y$.
Prove that every minimizer of $L_n$ has the unique form
\[
 w^\star+v,
 \qquad v\in\cR^\perp.
\]
Use the Pythagorean theorem to prove that $w^\star$ is the unique
minimum-Euclidean-norm minimizer.

\item Let $v_0$ be the orthogonal projection of $w_0$ onto $\cR^\perp$.
Prove that the $\cR^\perp$ component of every iterate equals $v_0$ and that
\[
 w_t\longrightarrow w^\star+v_0.
\]
Deduce the limit for $w_0=0$, and give an initialization whose limit is not
the minimum-norm minimizer.
\end{enumerate}
\begin{hints}
For (a), use the rank of $\mathbf X$ and the fact that a full-column-rank
matrix $\mathbf X^\top$ is one-to-one. For (b), first substitute $w^\star$
into $\mathbf Xw=y$. For (c), every increment belongs to $\cR$; on $\cR$, use
an eigenbasis of $\mathbf X^\top\mathbf X$.
\end{hints}
\end{exercisechunk}
